\section{Data and Measurement}\label{sec:data}

\subsection{Sample}

Daily open, high, low and close prices for S\&P~500 constituents were
obtained from Bloomberg. The panel keeps the 115 stocks with at least 70\%
of trading days observed between 29 November 2001 and 21 April 2026, a total
of 6136 trading days, and covers all eleven GICS sectors. Membership is that
of the April 2026 index, so the panel excludes firms that left the index
before then; Section~\ref{sec:discussion} discusses what this implies.
Returns are daily log price returns. They are winsorized point-in-time at the
expanding 0.1\% and 99.9\% quantiles after a 500-day warm-up, so that no
observation is trimmed using later data, and one special dividend that
appears as a price drop (Keurig Dr Pepper, 10 July 2018) is added back. The
market variables are the closing levels of the CBOE VIX index and the ICE
BofA MOVE index of Treasury implied volatility, and the three-month Treasury
bill rate.

\subsection{Variance}

Daily variance is measured by the range estimator of
\citet{parkinson1980extreme},
\begin{equation}
\sigma^2_{i,t} = \frac{\left(\ln H_{i,t} - \ln L_{i,t}\right)^2}{4\ln 2},
\label{eq:parkinson}
\end{equation}
where $H_{i,t}$ and $L_{i,t}$ are the day's high and low. The estimator uses
intraday information, is several times more efficient than the squared
close-to-close return, and is available for the full sample, which
five-minute realized variance is not. Days with a zero range are treated as
missing. All models work with the log of variance, which is close to
Gaussian and makes forecast errors comparable across stocks.

\subsection{Timing}

The forecast origin is the close of trading day $t$. Every predictor uses
information available at that moment: the day's range and return, and the
closing levels of the VIX and MOVE indices. The VIX is computed until 4:15
p.m.\ Eastern time, fifteen minutes after the equity close; we treat the two
as simultaneous. The forecast target for horizon $h$ is
\begin{equation}
y_{i,t,h} = \ln\Bigl(\frac{1}{h}\sum_{k=1}^{h} \sigma^2_{i,t+k}\Bigr),
\qquad h \in \{1, 5, 22\},
\label{eq:target}
\end{equation}
the log of mean variance over the next day, week and month. Forecasts are
made every fifth trading day, which gives 1078 weekly origins from 22
November 2004, the first date at which a 750-day window is available, to 21
April 2026. Section~\ref{sec:results} shows that this alignment matters.

\subsection{Full-sample memory}

Table~\ref{tab:memory} reports full-sample estimates of the fractional
parameter $d$ by the log-periodogram estimator of
\citet{geweke1983estimation} (GPH) and the local Whittle estimator of
\citet{robinson1995gaussian}, both with bandwidth $m=T^{0.65}$. Returns show
no memory: the mean estimate is $-0.02$. Log variance shows a great deal: the
mean GPH estimate is 0.54, the local Whittle estimate 0.52, and 84\% of stocks
have a GPH estimate above 0.5, the boundary of stationarity. These are the
familiar stylized facts \citep{andersen2003modeling}. A Hurst exponent of
0.06 from the scaling of increments of log variance would, read literally,
indicate very rough volatility \citep{gatheral2018volatility}; Appendix~\ref{app:rough}
shows that a two-factor Markovian model without roughness, observed through
the sampling noise of the range estimator, produces the same number, so
roughness is not identified at the daily frequency \citep{cont2024rough}. We do not use roughness further.

\begin{table}[H]
\caption{Full-sample memory estimates across the 115 stocks.}
\label{tab:memory}
\small
\begin{tabularx}{\textwidth}{Xcccccc}
\toprule
Series & Estimator & Mean & 10th pct & Median & 90th pct & Share $>0.5$ \\
\midrule
Returns & GPH & -0.020 & -0.082 & -0.022 & 0.044 & 0\% \\
 & LW & -0.027 & -0.077 & -0.027 & 0.025 & 0\% \\
Log Parkinson variance & GPH & 0.538 & 0.490 & 0.537 & 0.587 & 84\% \\
 & LW & 0.516 & 0.476 & 0.515 & 0.559 & 70\% \\
Log Parkinson variance & Hurst ($q=2$) & 0.060 & 0.053 & 0.060 & 0.069 & -- \\
\bottomrule
\end{tabularx}
\noindent{\footnotesize{\textit{Notes:} Daily data, 29 November 2001 to 21 April 2026 (6136 trading days). GPH: log-periodogram estimator of the fractional parameter $d$; LW: local Whittle; bandwidth $m=T^{0.65}$. Hurst: scaling of second moments of increments of log variance over lags 1 to 21 days. Percentiles are across stocks.}}
\end{table}


Estimates of $d$ above 0.5 for most stocks are themselves a warning. Long
memory and a slowly decaying short-memory process with occasional large
episodes are hard to tell apart in samples of this length
\citep{diebold2001long, granger2004occasional}. The rolling estimates that
follow make the distinction visible.
